\section{One Spiritual Family, Several Juridical Identities}

The Institute's public image is a family of priests, seminarians, oblates, Sister Adorers, and lay associates. Shared worship, patrons, symbols, and works make the language intelligible. Canonically and statistically, however, these groups are not interchangeable. The 2008 documents are especially valuable because they distinguish the male society from the women associated with it.

\begin{threestable}{Component}{Publicly documented identity}{Boundary}
Priests and seminarians & The society forms candidates and incorporates priestly members under proper law; its priests are styled canons. & A seminarian is not yet a priest or necessarily definitively incorporated.\\
Institute-designated clerical oblates & N. 181/2008 calls the early oblates lay; current pages describe non-priest full members dedicated to service. & The title does not make every oblate a canonical cleric; diaconate changes that status under canon 266.\\
Sister Adorers & N. 182/2008 erected a separate public association of faithful experimentally for three years and approved its constitutions for the same term. & The checked set contains no later act proving that the association became a society of apostolic life.\\
Society of the Sacred Heart & Its current institutional page presents a lay body sharing and supporting the spirituality. & Its members are not added to priest, seminarian, oblate, or Sister counts; a controlling canonical erection act was not located.\\
\end{threestable}

\subsection{Oblates: service without automatic priesthood}

Institutional vocation pages describe men who seek stable common life and service without a call to priesthood. Their work ranges across sacristy, music, administration, teaching, maintenance, art, and practical support. A multi-year formation introduces them to the Institute's languages, doctrine, history, liturgy, and spirituality. Public pages describe a perpetual obedience commitment and full membership in the Institute.

The phrase ``clerical oblate'' can mislead an outside reader. It is the Institute's current designation and reflects life within a priestly society. The 2008 decree explicitly calls the early group \emph{oblati laici}. Current pages allow that some may receive rites styled minor orders in Institute practice and potentially the subdiaconate or diaconate. Those rites do not by themselves create canonical clerical status under the current Code; ordination to the diaconate does. The publication therefore preserves the Institute's title while stating the legal boundary.

\subsection{The Sister Adorers and the decisive 2008 distinction}

The Sister Adorers of the Royal Heart of Jesus Christ Sovereign Priest describe a semi-contemplative and apostolic life centered on Eucharistic adoration, prayer for priests, common life, and assistance in Institute works. Their current site lists houses in several European countries and the United States and describes postulancy, novitiate, and profession.

Public institutional histories disagree on the beginning. The Sister site says the community was founded in January 2001; a United States ICKSP timeline likewise dates canonical establishment to 2001; an Italian legacy page says a community was canonically founded in 2004. No original act from either year was located to resolve whether the dates refer to first gathering, diocesan recognition, or another stage.

The Roman evidence in 2008 is exact. N. 181/2008 says the group was associated with the Institute but remained on the way toward the status of a society of apostolic life. The companion decree N. 182/2008 erected it as a public association of the faithful \emph{ad experimentum} for three years and approved its constitutions for three years. An older celebratory page calls the Sisters' act a pontifical-right decree; the dispositive text controls over that caption. No checked later status-changing act warrants calling the Sisters a pontifical-right society of apostolic life at the cutoff.

\subsection{Shared government is not presumed}

The Sisters' public site says they are integrated with the Institute and share its spirituality. The allegations reported in 2026 include claims about oversight relationships. Those facts do not allow this study to reconstruct the women's internal authority, vows, property, or the exact legal power of ICKSP superiors. Association, spiritual direction, and practical cooperation can be close while juridical persons remain distinct.

That distinction protects both accuracy and responsibility. Concerns raised by former Sister candidates must not automatically be attributed to every priest or treated as a finding against the male society. Conversely, calling the women ``separate'' cannot be used to deny documented shared identity, institutional links, or alleged involvement requiring competent review.

\statusframe{Pontifical Commission \emph{Ecclesia Dei}, N. 181/2008 and N. 182/2008, both 7 October 2008; current ICKSP, Sister, and Society of the Sacred Heart pages.}{The male society, its oblates, the women's group separately erected in 2008 as a public association, and the lay spiritual body belong to one public family but do not share one demonstrated canonical status.}{The full proper laws and any later Sister status act were not located. Current institutional terminology is reported without silently replacing the 2008 dispositions.}
